% Generated draft: requires mathematical and source review.
\documentclass[11pt]{article}
\usepackage[margin=25mm]{geometry}
\usepackage{fontspec}
\setmainfont{Latin Modern Roman}
\usepackage{amsmath,amssymb,mathtools,unicode-math}
\setmathfont{Latin Modern Math}
\usepackage{enumitem,needspace,xurl,hyperref}
\setlength{\parindent}{0pt}
\setlength{\parskip}{5pt}
\setlength{\emergencystretch}{3em}
\newcommand{\N}{\mathbb{N}}
\newcommand{\Z}{\mathbb{Z}}
\newcommand{\Q}{\mathbb{Q}}
\newcommand{\R}{\mathbb{R}}
\newcommand{\C}{\mathbb{C}}
\newcommand{\F}{\mathbb{F}}
\newcommand{\A}{\mathbb{A}}
\newcommand{\PP}{\mathbb{P}}
\newcommand{\E}{\mathbb{E}}
\newcommand{\eps}{\varepsilon}
\newcommand{\dd}{\,\mathrm{d}}
\begin{document}
\section*{UnsolvedMath ID: 30005147}
\textbf{Source problem code:} catalogue entry
\subsection{Definitions and mathematical statement}
A smooth manifold \(M\) admits positive scalar curvature if there exists a smooth Riemannian metric \(g\) on \(M\) such that
\[
\operatorname{Scal}_g(x)>0\qquad\text{for every }x\in M.
\]
This definition alone imposes neither completeness nor a boundary condition. The source uses the expression ``topologically PSC'' without defining it; whether it means admitting such a metric, or a more specific topological property, requires clarification.

In the usual smooth setting, a surgery of dimension \(k\) on an \(n\)-manifold is specified by a framed embedding \(S^k\times D^{n-k}\hookrightarrow\operatorname{int}(M)\). It removes \(S^k\times\operatorname{int}(D^{n-k})\) and attaches \(D^{k+1}\times S^{n-k-1}\) along their common boundary \(S^k\times S^{n-k-1}\). Here \(D^j\) denotes the closed \(j\)-dimensional disk and \(S^j\) the \(j\)-dimensional sphere. The surgery has codimension \(n-k\), so codimension at least \(3\) means
\[
0\leq k\leq n-3.
\]
This describes the standard surgery convention; the source does not explicitly specify the manifold category or surgery data.

\textbf{Source question.} Can every topologically positive-scalar-curvature \(n\)-manifold be built from ``simple'' topologically positive-scalar-curvature \(n\)-manifolds by performing surgeries of codimension at least \(3\)? The supplied assessment describes this as a question in all dimensions, rather than only dimension \(4\).

\textbf{Unresolved qualifications.} The source does not define ``simple,'' specify the range of \(n\), or state whether manifolds are smooth, closed, connected, oriented, compact, or allowed to have boundary. It also does not specify whether construction begins with a single simple manifold or a disjoint union of simple manifolds, whether only finitely many pieces and surgeries are permitted, or whether the resulting manifold must agree with the target up to diffeomorphism or homeomorphism. Consequently, the supplied record does not determine a fully quantified decomposition assertion. These choices must be supplied before the question has a precise mathematical formulation.
\subsection{Short English statement}
The question asks whether every manifold having the source's ``topologically PSC'' property can be assembled from simple manifolds with the same property using surgeries of codimension at least \(3\). The meaning of ``simple,'' the precise PSC property, and the permitted manifolds and constructions remain unspecified. The supplied curation labels the general question open and reports a four-dimensional analogue; these claims require independent review.
\subsection{Sources}
\begin{itemize}
\item Supplied local UnsolvedMath registry record, permanent ID \(30005147\), title ``Surgery Decomposition of Positive Scalar Curvature Manifolds,'' exact source problem code \texttt{catalogue entry}, in the Oberwolfach Reports Open Problems collection.
\item Supplied report citation: \emph{Geometrie (2023)}, Oberwolfach Reports, DOI \texttt{10.4171/owr/2022/28}, pp.\ 1551--1601. \url{https://doi.org/10.4171/owr/2022/28}. The record gives proposed year \(2022\).
\item The original excerpt attributes the contribution ``Decomposing 4-manifolds with positive scalar curvature'' to Christos Mantoulidis, jointly with Richard H. Bamler and Chao Li. No separate proposer is identified.
\item Supplied literature reference: Bamler--Li--Mantoulidis, ``PSC decomposition.'' \url{https://doi.org/10.1016/j.aim.2023.109231}.
\item Supplied literature reference: ``published article and stated Question 1.2.'' \url{https://www.sciencedirect.com/science/article/pii/S0001870823003742}.
\item Supplied literature reference: ``OWR report.'' \url{https://ems.press/content/serial-article-files/46965}.
\item Literature assessment checked at \texttt{2026-08-22}: the supplied curation labels the general all-dimensional surgery-decomposition question open and states that Bamler, Li, and Mantoulidis prove a substantial four-dimensional analogue, rather than the asserted statement for all \(n\). This is supplied dated curation derived from the OWR record, not independently verified here; its time-sensitive status and cited claims require independent verification.
\end{itemize}
\end{document}
